% Copyright (c) 2026 by C. Klukas.
% Licensed under the MIT License. See LICENSE for details.

A \paglet{} moves as a memory image. This chapter says what that image
holds and what it leaves behind, how a \paglet{} names its destination,
what the two hosts do between the call to \code{dispatch} and the
\code{arrived} event on the other side, and how clones, deactivation and
host restarts use the same image. Section~\ref{sec:lifecycle} gave the
events a \paglet{} sees around a move; here the view is the hosts'.

\subsection{What moves: the memory image}\label{sec:movement-image}

An instance of a WebAssembly module has three kinds of mutable state: its
linear memory, its mutable globals, and the stack of the function calls in
progress. A \paglet{} is the first two. The host takes an image only at a
quiescent point, between two handler calls, when no guest function runs
(Section~\ref{sec:model}); there is no call stack to save, and no native
stack, register or thread belongs to the \paglet{}. Everything the
\paglet{} keeps between messages lives in its linear memory: the
\code{Paglet} object with its members, the heap, the C++ shadow stack. The
host does not know what the bytes mean, and the \paglet{} writes no
serialization code: the journey's journal of Section~\ref{sec:lifecycle}
is a \code{std::vector} of structs in a member, and it arrives with the
\paglet{}.

% snippet: cpp/host/include/paglets/wasm/snapshot.hpp#image [float,caption={A memory image},label={lst:movement-image}]
\begin{lstlisting}[style=paglets-cpp,float,caption={A memory image \src{cpp/host/include/paglets/wasm/snapshot.hpp}},label={lst:movement-image}]
struct GlobalValue {
    std::string name;       // export name
    std::uint8_t type = 0;  // ValType byte
    std::uint64_t bits = 0;
};

struct PageEntry {
    bool zero = true;
    Digest hash{};  // unset for zero pages
};

struct Snapshot {
    Digest module_hash{};
    std::uint32_t page_count = 0;
    std::vector<GlobalValue> globals;
    std::vector<PageEntry> pages;    // page_count entries
    std::vector<std::uint8_t> data;  // non-zero pages, concatenated in order

    std::uint32_t data_pages() const { return static_cast<std::uint32_t>(data.size() / page_size); }
};
\end{lstlisting}
% /snippet

Listing~\ref{lst:movement-image} is the image. It names the module by its
SHA-256 hash, holds the number of 64\,KB pages of linear memory and the
value of every mutable global by export name, and for each page either the
mark ``zero'' or the page's SHA-256 hash with its content. The capture
(Listing~\ref{lst:movement-capture}) leaves zero pages out, the untouched
part of a \paglet{}'s memory, and hashes the rest. The hashes serve twice:
the receiving host checks every page against them, and it skips pages it
already holds (Section~\ref{sec:movement-protocol}).

% snippet: cpp/host/src/wasm_snapshot.cpp#capture [float,caption={Capturing an image},label={lst:movement-capture}]
\begin{lstlisting}[style=paglets-cpp,float,caption={Capturing an image \src{cpp/host/src/wasm\_snapshot.cpp}},label={lst:movement-capture}]
std::expected<Snapshot, std::string> capture(Instance& instance) {
    const ModuleInfo& info = instance.module()->info();
    if (auto ready = check_snapshot_ready(info); !ready) {
        return std::unexpected(ready.error());
    }

    Snapshot snap;
    snap.module_hash = instance.module()->hash();
    for (const Global& g : info.globals) {
        if (!g.is_mutable || g.imported) continue;
        auto bits = instance.global_bits(*g.export_name);
        if (!bits) return std::unexpected(bits.error());
        snap.globals.push_back({*g.export_name, static_cast<std::uint8_t>(g.type), *bits});
    }

    const std::span<std::uint8_t> mem = instance.memory();
    snap.page_count = static_cast<std::uint32_t>(mem.size() / page_size);
    snap.pages.resize(snap.page_count);
    for (std::uint32_t p = 0; p < snap.page_count; ++p) {
        const auto page = mem.subspan(std::size_t{p} * page_size, page_size);
        if (is_zero_page(page)) continue;
        snap.pages[p].zero = false;
        snap.pages[p].hash = sha256(page);
        snap.data.insert(snap.data.end(), page.begin(), page.end());
    }
    return snap;
}
\end{lstlisting}
% /snippet

Restoring an image instantiates the module, grows the memory to the
image's page count, copies the pages in and sets the globals
(\code{restore}, in the same file as the capture). The host does not call
\code{\_initialize} for a restore, but the engine does for every module
that imports WASI, which every guest built with the SDK does; the image
then overwrites whatever the initializer and the C++ static constructors
wrote. A side effect the initializer has through an import, such as a log
line, therefore happens again on every host the \paglet{} arrives on.
Only its first run, when the original \paglet{} was created, leaves a
trace in memory; clones start from the original's image.

\subsubsection{Globals}

The host cannot tell which mutable globals carry state a guest depends on
between handlers, and it does not try to. C and C++ guests compiled with
clang for \code{wasm32} keep their shadow stack pointer in a mutable
global, \code{\_\_stack\_pointer}, which is back at its initial value
between handlers; another toolchain may keep live state in a global. So
the rule is that every mutable global a module defines must be exported,
and \code{wasm-ld} does not export \code{\_\_stack\_pointer} unless asked.
The guest build (\code{cpp/cmake/PagletsGuest.cmake}) links every guest
with \code{-Wl,--export=\_\_stack\_pointer}, and the host checks the rule
statically when a module enters its store, before any instance exists
(Listing~\ref{lst:movement-snapshot-ready}).

% snippet: cpp/host/src/wasm_binary.cpp#snapshot_ready [float,caption={The check that a module can be imaged},label={lst:movement-snapshot-ready}]
\begin{lstlisting}[style=paglets-cpp,float,caption={The check that a module can be imaged \src{cpp/host/src/wasm\_binary.cpp}},label={lst:movement-snapshot-ready}]
std::expected<void, std::string> check_snapshot_ready(const ModuleInfo& info) {
    std::string missing;
    for (const auto& g : info.globals) {
        if (g.is_mutable && !g.imported && !g.export_name) {
            if (!missing.empty()) missing += ", ";
            missing += "#" + std::to_string(g.index);
        }
    }
    if (!missing.empty()) {
        return std::unexpected("mutable globals must be exported for memory images; not exported: " + missing +
                               " (link guests with -Wl,--export=__stack_pointer)");
    }
    return {};
}
\end{lstlisting}
% /snippet

One requirement the host cannot check: a module must not change its
function tables at run time, because function pointers in C++ are table
indices, and they stay meaningful across an image only while the table is
as the module defined it. Modules produced by clang satisfy this; the ABI
states it as a requirement (Appendix~\ref{app:abi}).

\subsubsection{Size and portability}

Linear memory is limited to 256 pages, 16\,MB
(\code{wasm::Limits::max\_memory\_pages}), and \code{restore} refuses an
image with more pages. Memory that a \paglet{} grew within the limit is
part of the image; its page count travels. The serialized form
(\code{PGIMG001}) is the same bytes on every host, and
\code{deserialize} checks every page hash before it returns. An image is
independent of CPU architecture and operating system because the module
is: the live mesh of Section~\ref{sec:evaluation} moves \paglets{}
between macOS on arm64 and Linux on aarch64 hosts without any step that
depends on either.

\subsection{What does not move}\label{sec:movement-record}

Besides the image, the host keeps a record for every \paglet{}: its ID,
module, trust class and owner, the capability table, the handle and
correlation counters, timers, requests in flight, residency marks and the
restart policy. Handles are host state: the \paglet{} holds an \code{i32},
the host holds what it stands for (Section~\ref{sec:model}). For a move
the record becomes the travelling state of
Listing~\ref{lst:movement-travel-state}, which the source host sends
beside the image.

% snippet: cpp/host/include/paglets/runtime/mobility.hpp#travel_state [float,caption={The state that travels beside the image},label={lst:movement-travel-state}]
\begin{lstlisting}[style=paglets-cpp,float,caption={The state that travels beside the image \src{cpp/host/include/paglets/runtime/mobility.hpp}},label={lst:movement-travel-state}]
// A paglet's state apart from its memory image, as it travels.
struct TravelState {
    PagletId id;
    std::string module;  // hex hash
    std::string trust;   // roaming, resident
    std::string owner;
    std::int32_t next_handle = 2;
    std::uint64_t next_correlation = 1;
    std::uint64_t next_timer = 1;
    std::int64_t checkpoint_ms = -1;
    std::vector<std::pair<std::string, std::int32_t>> services;   // default service endpoints: name, handle
    std::vector<std::pair<std::int32_t, Cap>> caps;               // every capability, endpoints with host hints
    std::vector<std::pair<std::uint64_t, std::int64_t>> pending;  // requests in flight: correlation, ms left
    std::vector<std::string> marks;                               // data residency marks
    bool restart = true;                                          // restart policy (abi::RestartArg)
    bool notify_recovered = false;
};
\end{lstlisting}
% /snippet

A capability may stand for a resource of the host it was made on. The
destination therefore rebuilds the capability table entry by entry (the
runtime's \code{prepare\_arrival}); Table~\ref{tab:movement-caps} says
what becomes of each kind. The handle
numbers stay the same, so handles the \paglet{} stored in its memory
remain valid. An entry the destination cannot rebuild is dropped, and the
\code{arrived} event lists it in \code{lost}, as
\code{handle 7: dir data/docs}, so the \paglet{} can ask for it again or
give up on that work.

\begin{table}[htbp]
\centering\small
\caption{Capabilities on arrival (\code{Runtime::prepare\_arrival}, \code{Node::Impl::recreate}).}
\label{tab:movement-caps}
\begin{tabular}{p{0.40\linewidth}p{0.52\linewidth}}
\toprule
Capability & On the destination \\
\midrule
Endpoint to a \paglet{} & kept; it names the host its target is on \\
Default endpoint to a system service & replaced by the destination's endpoint to the same service, under the same handle; lost when the destination does not run the service \\
Reply capability & kept \\
Timer & re-armed for its time \\
Resource from a grant (a \code{dir} of the \code{files} service, an endpoint to a system service that a grant gave) & re-created from the grant when the grant is the \paglet{}'s own, not expired and covers the destination; narrowings the \paglet{} made stay \\
Pin & kept: it refers to the pinned \paglet{}, and every host has a locator \\
Artifact & kept: it names the host that holds the content \\
Anything else & lost; listed in \code{ArrivedEvent::lost} \\
\bottomrule
\end{tabular}
\end{table}

Two host resources are outside the table. The key-value store of the
\code{storage} service is a directory of the host; the runtime removes it,
together with the \paglet{}'s scratch directory, when a move commits
(\code{finish\_departure} in \code{cpp/host/src/runtime.cpp}). A \paglet{}
that needs data on the next host carries it in memory or shares it as an
artifact (Section~\ref{sec:mesh}). The compiled form of the module does
not move either: every host translates the module for its own engine.

\subsection{Transfer tickets}\label{sec:movement-tickets}

\code{paglets::dispatch(destination)} takes a transfer ticket, a string
of the form
\begin{center}
\code{TARGET[?retries=N\&arrival=active|inactive]}
\end{center}
(Table~\ref{tab:movement-tickets}). The call returns at once; the move
begins when the handler returns. It returns an error without moving
anything when the host is not in a mesh (\code{unsupported}), when the
\paglet{} is a system or resident \paglet{} (\code{denied}; only roaming
\paglets{} move), when another lifecycle operation is already pending in
this handler (\code{bad\_state}), or while the \paglet{} is pinned
(\code{pinned}, Section~\ref{sec:communication}).

\begin{table}[htbp]
\centering\small
\caption{Transfer tickets (\code{parse\_ticket} and \code{candidates} in \code{cpp/host/src/node/movement.cpp}).}
\label{tab:movement-tickets}
\begin{tabular}{p{0.30\linewidth}p{0.62\linewidth}}
\toprule
Target & Hosts it names \\
\midrule
a host name, a key ID, or a prefix of a key ID of at least 8 hex digits & that host \\
\code{label:L} & every host with label \code{L}, in random order \\
\code{offer:S} or \code{offer:S.OP} & the hosts that offer service \code{S} (operation \code{OP}), the least loaded first \\
\code{any} & every other host, in random order \\
\midrule
\code{?retries=N} & try at most \code{N} of the matching hosts, 1 to 16 (default 3) \\
\code{?arrival=inactive} & store the \paglet{} on arrival; it runs when its next message comes \\
\bottomrule
\end{tabular}
\end{table}

The source host turns the ticket into a list of candidates. It takes the
enrolled hosts other than itself that match the target and leaves out
hosts that run another mesh protocol or paglet ABI, or an older minor ABI
version (the \paglet{} may use newer imports). It then applies the mesh's
rules: a host where an item of the \paglet{}'s manifest would be denied is
skipped unless a grant for that host covers it, and a host that a
residency rule excludes is skipped (Section~\ref{sec:movement-code}).
Hosts that are up come first, and the list is cut to \code{retries}. Every
host left out is remembered with its reason; when no candidate remains,
the \paglet{} sees them all: \code{no host matches label:gpu; host lab-2:
the manifest's files access is denied there}. When candidates remain and
every one refuses, the failure names the refusals
(Section~\ref{sec:movement-protocol}).

The journey's \code{go\_on} (Listing~\ref{lst:lifecycle-journey-go-on}
in Section~\ref{sec:lifecycle-moves}) is its step from one host to the
next. The itinerary is a vector of key IDs that the \paglet{} filled
from the landscape of \code{mesh-info} on its home host; \code{next\_},
\code{home\_} and \code{left\_ms\_} are members like any other, and they
are the same values on the next host. \code{dispatch} fails at once only
for an empty destination or a \paglet{} that may not move (the cases
listed above). A malformed ticket
(\code{retries=0}, an unknown parameter) is found when the move starts,
after the handler returned, and reported like a destination that no host
matches: through \code{on\_move\_failed} on the old host, while a
successful move ends with \code{on\_arrived} on the new one.

\subsection{The move protocol}\label{sec:movement-protocol}

A move is a conversation between two hosts, with the hosts that keep the
\paglet{}'s location record as a third party
(Figure~\ref{fig:movement-protocol}). The code is
\code{cpp/host/src/node/movement.cpp} on both sides; the runtime's part is
\code{Runtime::finish\_departure}, \code{prepare\_arrival} and
\code{commit\_arrival}.

\begin{figure}[htbp]
\centering
\begin{tikzpicture}[x=1cm,y=1cm]
  \node[box] (S) at (0,0) {source host};
  \node[box] (D) at (7.4,0) {destination host};
  \node[box] (L) at (13.0,0) {location holders};
  \draw[lifeline] (S.south) -- ++(0,-9.4);
  \draw[lifeline] (D.south) -- ++(0,-9.4);
  \draw[lifeline] (L.south) -- ++(0,-9.4);
  \node[event,anchor=west,align=left] at (0.15,-0.75) {\code{dispatching} event, image taken, messages held};
  \draw[msg] (0,-1.4) -- node[above]{\code{move-offer}: signed envelope, state, page manifest} (7.4,-1.4);
  \node[event,anchor=west,align=left,text width=4.9cm] at (7.55,-2.05) {checks: enrolled sender, signature, passport, module trust, not here yet};
  \draw[dashedarrow] (7.4,-2.75) -- node[above,font=\footnotesize,color=pagletsMuted]{\code{module-want} / \code{module-part}, if the module is missing} (0,-2.75);
  \draw[msg] (7.4,-3.45) -- node[above]{\code{move-want}: the pages not in its cache} (0,-3.45);
  \draw[msg] (0,-4.15) -- node[above]{\code{move-pages}: compressed, checked against the manifest} (7.4,-4.15);
  \node[event,anchor=west,align=left,text width=4.9cm] at (7.55,-4.8) {prepared: owner quota, capabilities rebuilt, not yet running};
  \draw[msg] (7.4,-5.5) -- node[above]{\code{move-ready}} (0,-5.5);
  \draw[msg] (0,-6.2) -- node[above,pos=0.8]{\code{loc-set}: the new record} (13.0,-6.2);
  \draw[msg] (13.0,-6.9) -- node[above,pos=0.26]{\code{loc-ack} from a majority} (0,-6.9);
  \node[event,anchor=west,align=left] at (0.15,-7.5) {decision on disk; held messages follow};
  \draw[msg] (0,-8.2) -- node[above]{\code{move-commit}} (7.4,-8.2);
  \node[event,anchor=west,align=left,text width=4.9cm] at (7.55,-8.85) {stored, then started with \code{arrived}};
  \node[event,anchor=north west,align=left,text width=13.6cm] at (-0.6,-10.0) {At any step: \code{move-refused}, or no answer within 30\,s (\code{move-abort}), and the source offers the \paglet{} to the next candidate. A prepared destination that hears nothing asks \code{move-query} until the source answers.};
\end{tikzpicture}
\caption{A move between two hosts. Solid arrows are frames on the encrypted host channels; the dashed one happens only when the destination lacks the module.}
\label{fig:movement-protocol}
\end{figure}

\subsubsection{Departure}

When the handler that called \code{dispatch} returns, the runtime
delivers the \code{dispatching} event, takes the image, and checks that
no pin arrived meanwhile. It then releases the instance and holds every
message that arrives for the \paglet{}; until the mesh reports the
outcome the \paglet{} exists only as its image and record. The node
attaches the \paglet{}'s passport, the signed chain from an enrolled
owner that says who launched it (Section~\ref{sec:mesh}); a \paglet{}
without one does not move (\code{no passport: only paglets of enrolled
owners move}). Before it looks at candidates, the node applies the
residency rule that forbids leaving at all
(Listing~\ref{lst:movement-residency}).

% snippet: cpp/host/src/node/movement.cpp#residency [float,caption={Data that stays on its host keeps the \paglet{} there},label={lst:movement-residency}]
\begin{lstlisting}[style=paglets-cpp,float,caption={Data that stays on its host keeps the \paglet{} there \src{cpp/host/src/node/movement.cpp}},label={lst:movement-residency}]
// Data residency (planning/cpp-residency.md): a paglet that read a root
// whose content stays on this host does not leave.
for (const auto& mark : d.state.marks) {
    if (!mark.starts_with("root:")) continue;
    auto rule = ledger.state().residency.find(mark.substr(5));
    if (rule != ledger.state().residency.end() && rule->second.host_only)
        return fail("data residency: the content of root " + rule->first + " stays on this host");
}
\end{lstlisting}
% /snippet

\subsubsection{The offer and the checks}

The offer is a canonical map: protocol version, mesh, source and
destination keys, request number, kind (\code{dispatch} or \code{clone}),
\paglet{} ID, passport, travelling state, the page manifest (page count,
globals, and per page a hash or nil), the arrival mode, the move counter
after this move, and the time. The source signs the SHA-256 of the body
with its host key under the domain string \code{paglets move offer v1}, so
the destination can check who sent it independent of the channel the
frame came through, which matters for relayed hosts. Before any page
moves, the destination refuses unless everything in
Listing~\ref{lst:movement-accept} holds: the sender is an enrolled host
and the signature verifies, the offer names this host and this mesh, the
state decodes and agrees with the offer, the passport verifies for this
\paglet{} and module (its owner is enrolled and the passport has not
expired; for clones and children, every link is signed by an enrolled
host), the module may run as a roaming \paglet{} in this mesh, and the
\paglet{} is not here already.

% snippet: cpp/host/src/node/movement.cpp#accept [float,caption={What the destination checks before any page moves},label={lst:movement-accept}]
\begin{lstlisting}[style=paglets-cpp,float,caption={What the destination checks before any page moves \src{cpp/host/src/node/movement.cpp}},label={lst:movement-accept}]
if (*mesh_id != ledger.mesh() || *body_from != from || *to != key.public_key() ||
    static_cast<std::uint64_t>(*body_request) != request || (*kind != "dispatch" && *kind != "clone")) {
    return refuse("offer not for this host");
}
auto state = runtime::decode_state(*state_bytes);
if (!state) return refuse(state.error());
auto module = mesh::parse_key_id(state->module);
if (!module || state->id != *paglet) return refuse("offer and state disagree");
auto passport = mesh::Passport::decode(*passport_bytes);
if (!passport) return refuse("passport: " + passport.error());
if (auto ok = mesh::verify_passport(*passport, st, *paglet, *module, mesh::unix_ms()); !ok) {
    return refuse("passport: " + ok.error());
}
if (auto verdict = mesh::check_module(st, *module, "roaming"); !verdict.allowed) return refuse(verdict.reason);
if (runtime.info(*paglet)) return refuse("paglet " + paglet->substr(0, 8) + " is here already");
auto image = parse_image_manifest(*manifest, *module);
if (!image) return refuse("malformed page manifest");
\end{lstlisting}
% /snippet

\subsubsection{The module and the pages}

A destination that lacks the module fetches it by hash: from its module
sources (directories the host was started with), then from the source
host, then from every other host of the mesh, in parts of at most 1\,MB
(\code{need\_module} in \code{cpp/host/src/node/node\_impl.hpp}). Whatever
arrives is verified against the hash and checked as a \paglet{} module
before it is stored; a host that sends something else or does not answer
within 10\,s is skipped for the next. Modules are at most 64\,MB. A host
that still holds the module, because it ran it before or received it with
an earlier move, fetches nothing; a module without a \paglet{} is kept for
24 hours (Section~\ref{sec:movement-code}).

Pages are transferred only when needed. Every host keeps the pages of
images it sent or received in a cache of 256\,MB, least recently used
out, keyed by page hash. The destination compares the manifest against
the cache and asks (\code{move-want}) only for the page indices it lacks;
the source answers with those pages, each compressed with zstd, in frames
of about 1\,MB (a frame is closed once its compressed pages reach that
size). A \paglet{} that returns to a host it visited, or whose
memory changed little since its last move, transfers only the pages that
changed. Every page received is decompressed to exactly 64\,KB and hashed;
a page that does not match the signed manifest ends the arrival
(\code{damaged page}).

With all pages present, the destination prepares the arrival: the runtime
creates the \paglet{} without running it, checks the module's admission
as a roaming \paglet{} and the owner's quota on this host (10\,000
\paglets{} of one owner per host), rebuilds the capability table
(Table~\ref{tab:movement-caps}), re-arms timers and the timeouts of
requests in flight, and keeps the \code{arrived} event ready. Then it
answers \code{move-ready}.

\subsubsection{Commit}

Only the source decides. On \code{move-ready} it first sends the
\paglet{}'s new location record, with the incremented move counter, to the
hosts responsible for it (three of the hosts that are up, chosen by
consistent hashing, or all of them when three or fewer are up,
Section~\ref{sec:communication}) and waits for a majority of
acknowledgements. Then it writes the decision to its state directory,
tells the runtime, and sends \code{move-commit}. The runtime forwards the
held messages to the new host, remembers for ten minutes where the
\paglet{} went (a tombstone, so endpoints with a stale location keep
working), removes the \paglet{}'s record, image and directories, and
reports the outcome to anyone waiting. The destination, on
\code{move-commit}, writes the \paglet{} and its passport to disk at once
and starts it with \code{arrived}, or keeps it stored until a message
comes when the ticket said \code{arrival=inactive}.

The order protects the invariant that a \paglet{} runs on at most one
host. The destination never starts a prepared \paglet{} on its own: when
it hears nothing within the move timeout of 30\,s it asks
\code{move-query}, and keeps asking, and the source answers from its
recorded decisions (kept for an hour) with \code{move-commit} or
\code{move-abort}. A source that crashed after writing the decision
answers the same way when it comes back. If the location records cannot
be updated within the move timeout, the source aborts the move. A
responsible host that answers with a newer record (a counter lost in a
crash) makes the source record again above it, at most three times, and
then abort as well. An aborted move overwrites the records with a higher
counter that names the source: the \paglet{} stays.

\subsubsection{Failure}

The destination may refuse, with the reason (\code{move-refused}), or
stay silent for 30\,s, after which the source sends \code{move-abort};
either moves the offer to the next candidate of the ticket, each offer
with a new request number so late answers do not match. When no candidate is left the move fails; the
runtime makes the \paglet{} ready again with its held messages and
delivers \code{paglets.move\_failed} with the destination it asked for and
the reason of every host tried, such as \code{host lab-2: quota: its
owner has as many paglets here as one owner may have}. Nothing of the
\paglet{} is lost in any of these cases, because the image on the source
is the \paglet{} until the commit.

\subsubsection{From the command line}

An owner or admin moves a \paglet{} with \code{paglets-host remote
dispatch}, which sends the ticket to the host the \paglet{} is on and
waits for the outcome through the same hook the runtime calls at the end
of a move (Listing~\ref{lst:movement-cli}; the session options are those
of Section~\ref{sec:tools-cli}):

\begin{lstlisting}[style=paglets-bash,caption={Moving a \paglet{} from the command line (\code{paglets-host remote}).},label={lst:movement-cli}]
H=paglets-host
S="--connect https://lab-1:7443 --key olga.key --ledger ledger"
ID=$($H remote launch $S build/macos-arm64/guests/journey.wasm)
$H remote dispatch $S $ID lab-2            # moved to lab-2 2f8e1c0b94a7d3e1
$H remote dispatch $S $ID "label:gpu?retries=2"
# the move failed, the paglet stays: no host matches label:gpu
$H remote locate $S $ID                    # host:  lab-2 2f8e1c0b94a7d3e1...
                                           # moves: 1
\end{lstlisting}

\code{remote dispatch} waits at most 120\,s for the outcome and exits
with status 1 and the reason when the move failed.
\code{paglets-web} offers the same move (Section~\ref{sec:tools}).

\subsection{Clones on other hosts}\label{sec:movement-clones}

\code{clone\_raw(args, caps, destination)} makes a copy of the \paglet{}
with the same image, and with a destination the copy is made on another
host. The image is the one taken when the handler returns, so all clones
of one handler call share it; at most 16 children and clones may be made
per call. A remote clone leaves as a departure of its own, with the
clone's ID in the travelling state. It keeps the original's transferable
endpoints and its default service endpoints under the same handle numbers
(so handles in the copied memory stay meaningful), but no requests in
flight; \code{caps} may hold only endpoints to \paglets{} and artifacts,
the capabilities that are meaningful on every host
(\code{invalid\_argument} otherwise). The original continues;
\code{clone\_raw} returns an endpoint
to the clone at once, which reaches it wherever it lands. On the
destination the clone goes through the same offer, checks and commit as a
move, with a passport link that the source host signed to extend the
original's passport, and starts with \code{cloned} (the arguments, the
capabilities and the original's ID) instead of \code{arrived}. If no host
takes it, the original receives \code{on\_move\_failed} with the clone's
ID in \code{MoveFailed::clone}.

The inventory example uses this for a fan-out: one clone per host, each
asks its host's \code{server-info} and reports to the parent. The parent
starts the fan-out with the hosts from \code{mesh-info}'s landscape and
passes a failed clone to the pattern (Listing~\ref{lst:movement-fanout});
the \code{FanOut} pattern of the SDK's patterns library
(\code{patterns/fanout.hpp}) clones the \paglet{} to each destination with
the parent's own endpoint as
the only capability, and arms a timer for the deadline.

% snippet: cpp/examples/inventory/inventory.cpp#fanout [float,caption={The inventory's fan-out},label={lst:movement-fanout}]
\begin{lstlisting}[style=paglets-cpp,float,caption={The inventory's fan-out \src{cpp/examples/inventory/inventory.cpp}},label={lst:movement-fanout}]
    void on_move_failed(const paglets::abi::MoveFailed& f) override {
        if (!f.clone.empty()) fan_.clone_failed(f);
    }

private:
    void start(const Collect& c, Message& m) {
        report_ = Report{};
        auto reply = std::make_shared<paglets::Capability>(m.defer_reply());
        auto go = [this, c, reply](std::vector<std::string> hosts) {
            const auto sent = fan_.start(hosts, paglets::encode(c), c.timeout_ms);
            (void)paglets::reply_to(*reply, Started{sent > 0, sent > 0 ? "" : "no clone went out",
                                                    static_cast<std::int64_t>(hosts.size())});
        };
        if (!c.hosts.empty()) return go(c.hosts);
        auto mi = paglets::service("mesh-info");
        if (!mi) return go({""});
        svc::mesh_info::Client client{*mi};
        auto sent = client.landscape({}, [go](paglets::Result<svc::mesh_info::Landscape> l, Message&) {
            std::vector<std::string> hosts{""};
            if (l) {
                for (std::size_t i = 1; i < l->hosts.size(); ++i) hosts.push_back(l->hosts[i].host);
            }
            go(std::move(hosts));
        });
        if (!sent) go({""});
    }
\end{lstlisting}
% /snippet

A clone finds itself on the new host in \code{on\_cloned}
(Listing~\ref{lst:movement-clone}). Its first capability is the parent's
endpoint; it looks up \code{server-info} on the host it is on, asks it,
reports through the endpoint and disposes itself.

% snippet: cpp/examples/inventory/inventory.cpp#clone [float,caption={A clone on its host},label={lst:movement-clone}]
\begin{lstlisting}[style=paglets-cpp,float,caption={A clone on its host \src{cpp/examples/inventory/inventory.cpp}},label={lst:movement-clone}]
void on_cloned(paglets::Started& s) override {
    Collect q;
    if (s.caps.empty() || !s.decode_args(q)) {
        (void)paglets::dispose();
        return;
    }
    parent_ = paglets::Endpoint(s.caps[0].release());
    query_ = q;
    pt::lookup("server-info", [this](paglets::Result<paglets::Endpoint> ep) {
        if (!ep) return finish(pt::error_text("server-info", ep.error()));
        info_ = *ep;
        summary();
    });
}
\end{lstlisting}
% /snippet

The parent sees each child once, when it reported, when its clone could
not be made or did not arrive, or when the deadline passed without a
report (Listing~\ref{lst:movement-collect}). An inventory is therefore
complete even when a host is down: that host's entry carries the error.

% snippet: cpp/examples/inventory/inventory.cpp#collect [float,caption={Collecting the reports},label={lst:movement-collect}]
\begin{lstlisting}[style=paglets-cpp,float,caption={Collecting the reports \src{cpp/examples/inventory/inventory.cpp}},label={lst:movement-collect}]
fan_.attach(router(), [this](const pt::FanOut::Child& c) {
    HostInventory h;
    if (!c.ok || !paglets::decode(c.result, h)) {
        h = HostInventory{};
        h.error = c.error.empty() ? "no report" : c.error;
    }
    h.host = c.host;
    h.host_name = c.host_name.empty() ? c.destination : c.host_name;
    report_.hosts.push_back(std::move(h));
    report_.finished = fan_.finished();
    if (report_.finished) {
        std::ranges::sort(report_.hosts, [](const auto& a, const auto& b) { return a.host_name < b.host_name; });
    }
});
\end{lstlisting}
% /snippet

\subsection{Deactivation, checkpoints and restarts}\label{sec:movement-disk}

The image that moves a \paglet{} also parks it. \code{deactivate()} takes
effect when the handler returns: the \code{deactivating} event is
delivered, the image is taken, and the instance is released, so an
inactive \paglet{} holds no memory and no compiled code. On a host with a
state directory the image is written to \code{paglets/<id>.paglet}
together with the record, to a temporary file first and then renamed, so
a crash leaves the old file or the new one; without a state directory the
image stays in the host's memory. The next message, timer or wake-up
restores the instance from the image and delivers \code{activated}
before it (Section~\ref{sec:lifecycle-deactivation}).

The same file is the \paglet{}'s checkpoint. With a state directory, the
runtime stores an image after a handler when the last one is older than
the checkpoint interval (1\,s by default, per \paglet{}), at once after
\code{checkpoint()}, and at once when a moved or cloned \paglet{} is
committed, so a crash right after a commit finds the \paglet{} on the
destination. The restart policy of Section~\ref{sec:lifecycle-restart}
(\code{set\_restart(false)}) turns this off.

When a host starts, it reads every \code{.paglet} file and re-creates the
\paglets{} with the restart policy: the record with its capabilities,
timers re-armed for their time, and the image to be restored on the first
delivery (\code{Runtime::Impl::recover} in \code{cpp/host/src/runtime.cpp}).
Requests that were in flight at the crash are answered \code{timeout}, so
the handlers waiting for them run and can retry. A \paglet{} that asked
for it receives \code{recovered} first (Section~\ref{sec:lifecycle}). A
\paglet{} whose module is no longer in the store is dropped with a log
line. A crashed worker process takes the same path for its \paglets{}:
they resume from their last image, except the one whose handler was
running.

\subsection{Code mobility and its limits}\label{sec:movement-code}

Code is named by content. A \paglet{}'s record names its module by SHA-256
hash, the image is bound to that hash, and every host that runs the
\paglet{} holds the module under that hash in its store, verified when it
arrived and again whenever it is read from disk. A module that a host
fetched for one arrival serves every later one; modules without a
\paglet{} are collected after 24 hours. Hosts send modules only to
enrolled hosts.

Whether a module may run is a decision of the mesh, recorded in the
ledger (Section~\ref{sec:mesh}). Signers vouch for a module with a
signature record; admins trust a signer, or a module hash, for the
classes roaming, resident and system; revoking the signer's key, the
signature or the module hash takes the trust back. By default roaming
modules need no trust (\code{module-policy any}); with
\code{module-policy trusted} the destination refuses a module that no
trust record covers before any page moves, and a host ends the
\paglets{} of a module that loses trust while they run. An arriving \paglet{} is always checked as
roaming, whatever it was: resident and system \paglets{} do not move.

Data residency follows the data into the \paglet{}. When a \paglet{} reads
from a named root through the \code{files} service, its host marks it
\code{root:NAME}; the mark is never removed, travels with the \paglet{}
and passes to its children and clones. A residency rule in the ledger
makes the mark bind: \code{host-only} keeps the \paglet{} where it is
(Listing~\ref{lst:movement-residency}: the host where it read, unless
the rule came after a later move), \code{hosts}
limits the candidates of every ticket to the hosts a selector names, by
key or label. A rule made after the read applies, because the mark names
the root, not the rule. The content exchange between hosts obeys the
same rules for artifacts and shared files.

The design has these limits:

\begin{itemize}
\item Only roaming \paglets{} of enrolled owners move, only within one
  mesh, and only to hosts that run the same major and at least the same
  minor ABI version.
\item A \paglet{} never changes code while it lives. The image belongs to
  one module hash; a new version is a new module, and its \paglets{}
  start fresh. There is no state migration between versions.
\item Moves happen between handlers. A computation that takes longer than
  one handler (5\,s) must be split into messages the \paglet{} sends to
  itself, and it can move only between them. This is weak mobility in the
  sense of Section~\ref{sec:background}, by design.
\item A \paglet{} is at most 16\,MB of memory, images included. Larger
  data stays on hosts and travels as artifacts.
\item Capabilities that stand for host resources are re-created only from
  grants; the \code{storage} store and the scratch directory do not
  follow the \paglet{}.
\item The host cannot verify that a module leaves its tables alone; a
  module that does not satisfies the ABI only on paper.
\end{itemize}

Section~\ref{sec:evaluation} measures what these decisions cost: image
sizes, the time to capture and restore an image, the pages a repeated
move sends, moves between two hosts in one process, and the transfer
rates between the hosts of the live mesh. A move between two real
machines is exercised there but not timed.
